\scanpage{15}
\mspage{PU 37}
\deleted{traditionnellement, de paire d' o p p o s é s , comme c'est le cas pour les couples voisins
\[
\text{le précis }\handwritten{\longrightarrow}\text{ le vague}\qquad\text{ou}\qquad
\text{le particulier }\handwritten{\longrightarrow}\text{ le général , -\unsure{.}}
\]
et p\deleted{a}\inserted{\handwritten{\unsure{ou}}}r les quatre autres couples qui suivent dans notre zig-zag. Je présume que même de nos jours, il y aura peu de scientifiques qui songeraient à récuser ce couple, comme exprimant l'harmonie idéale recherchée dans la démarche scientifique.

Pour ma part, si quelqu'un s'était avisé de me poser la question, nul doute que je me serais reconnu dans cet idéal-là, sans y réfléchir à deux fois. Aujourd'hui encore je ne le renierais pas. Mais, comme tout idéal, celui-ci ne touche que la pure surface des choses. Ce n'est pas dans l'idéal, mais dans la riche r é a l i t é et dans sa mère, le r ê v e , que réside la profondeur, et la fécondité qui est sienne.

13.\quad Les épousailles de l'ordre et du mystère

(23 mars)\quad La nuit dernière, en suivant les associations de couples yin-yang suscitées par la réflexion, j'ai vu se prolonger encore sur la gauche notre joli zig-zag de hier, de façon à entrer dans des types de qualités qui ne sont pas particulières seulement à la pensée. Dans le diagramme initial, j'ai remplacé les termes ``le général'' (donc des adjectifs ``substantivisés'') par les noms correspondants ``généralités'', etc. Cela conduit alors à le prolonger comme suit (*) :

\diagram{\begin{tikzcd}[ampersand replacement=\&,column sep={6.7mm,between origins},row sep=10pt,cells={nodes={font=\tiny,inner sep=0.5pt}}]
{} \& \text{discipline}\handwritten{\arrow[dl]\arrow[dr,line width=.52pt]} \& {} \& \text{contrôle}\handwritten{\arrow[dl]\arrow[dr,line width=.52pt]} \& {} \& \text{volonté}\handwritten{\arrow[dl]\arrow[dr,line width=.52pt]} \& {} \& \text{rigueur}\handwritten{\arrow[dl]\arrow[dr,line width=.52pt]} \& {} \& \text{précision}\handwritten{\arrow[dl]\arrow[dr,line width=.52pt]} \& {} \& \text{clarté}\handwritten{\arrow[dl]\arrow[dr,line width=.52pt]} \& {} \& \text{savoir}\handwritten{\arrow[dl]\arrow[dr,line width=.52pt]} \& {}\\
\text{imagination} \& {} \& \text{jeu} \& {} \& \text{abandon} \& {} \& \text{spontané\deleted{\unsure{te}}\handwritten{\unsure{it}}é} \& {} \& \text{généralité} \& {} \& \text{flou} \& {} \& \text{obscurité} \& {} \& \text{mystère ,}
\end{tikzcd}}

où la ligne du haut est encore formée des termes yang, celle du bas des termes yin.

\noindent\rule{2.1cm}{0.4pt}

(*) (25 mars)\quad Avant de prolonger le zig-zag de la veille, j'avais été conduit à remplacer le terme yang ``le particulier'' (s'accouplant avec ``le général'' ou ``généralité'', et faisant presque double-emploi avec son voisin de droite ``précision'') par le terme ``rigueur'', pour former le nouveau couple
\[
 \text{rigueur - généralité}
\]
venant parfaire son voisin (ou ``concubin'') de droite
\[
 \text{précision - généralité .}
\]
C'est bien le premier de ces deux couples qui me paraît le mieux ``assorti''. S'il est vrai que souvent, on perd en précision ce qu'on gagne en généralité et inversement, une telle situation ne se présente j a m a i s pour le couple ``rigueur -}
% archive sheet number: 14

